% ---------------------------------------------------------------------
% Drop-in continuation of Sec. VI-B (ground-effect-induced bias) for the
% IEEE Access manuscript.  Place after the LPV affine fit of the channel
% gains (Table 2 / Fig. 15); it takes the coefficients from there and
% adds what the excitation data say about them.
%
% Cross-references assumed to exist in the host document:
%   \eqref{eq:mcrit_roll}   Eq. (7)    directional roll thresholds
%   \eqref{eq:mcrit_pitch}  Eq. (14)   directional pitch thresholds
%   \eqref{eq:weight}       Eq. (34)   weight from the moments
%   \eqref{eq:offset}       Eq. (35)   offset per direction
%   \eqref{eq:pivotfree}    Eq. (36)   pivot-free average
% Replace these with the host's own labels before compiling.
% ---------------------------------------------------------------------

\subsubsection{What the excitation data say about the bias}

The coefficients above are predictions. The excitation runs test them
without any quantity being fitted, through the gap that the rigid
balance leaves. Writing the rigid threshold of
\eqref{eq:mcrit_roll} and \eqref{eq:mcrit_pitch} as
$T_{\mathrm{rigid}} = \mathrm{sgn}\,(W - f)\,l + S_{\mathrm{off}}$ and
the detected onset moment as $M_{\mathrm{det}}$,
\begin{equation}
\Delta M \;\triangleq\; T_{\mathrm{rigid}} - M_{\mathrm{det}}
\label{eq:gap}
\end{equation}
is the moment the ground supplies but the rigid balance omits. Over
the twenty direction groups of the five configurations, the vehicle
tips at a \emph{smaller} commanded moment than the rigid balance
requires in every one of them, by $204$~mN\,m in the median --- about
a fifth of the threshold itself. Evaluated at the onset, the
interference model of \eqref{eq:ge_moment} predicts $165$~mN\,m of
that, or $81\%$, with no tuned constant; the single-rotor reference,
which neglects rotor--rotor interference, predicts a fifth of it.
Group by group the ratio of implied to modelled spans $0.53$ to
$2.07$, so this is a statement about the median and not about any one
group.

One qualification belongs with the number. $\Delta M$ is not a
measurement of the aerodynamic moment: it is what the rigid balance
leaves over, and two of that balance's ingredients are themselves
modelled --- the collective thrust is $\sum C_T \omega_i^2$ from a
bench identification, hence an out-of-ground-effect value by
construction, and the circle fit of Sec.~V locates the
\emph{kinematic} pivot, whereas the balance needs the line of action
of the normal-force resultant. $\Delta M$ therefore collects the
ground's total contribution, aerodynamic and contact alike.

\paragraph{Two checks that use no external reference}
Both \eqref{eq:weight} and \eqref{eq:offset} close on themselves.
Eq.~\eqref{eq:offset} estimates the same offset twice, once from each
tip direction with that direction's own arm, and \eqref{eq:weight}
returns the weight from the moments. Neither closes without a
ground-effect term: the two directional offset estimates disagree by
$12.6$~mm in the median ($19.8$ at worst), and \eqref{eq:weight}
returns a weight $5.2\%$ below $mg$, low in all ten case/axis
configurations. Restoring the channels brings the directional spread
to $3.6$~mm ($9.0$ at worst) and the weight to $1.1\%$ below; the
thrust channel does most of the latter, since it enters
\eqref{eq:weight} as $(1+\eta_f^{GE})f_{\mathrm{crit}}$. Neither the
load cell nor the identified offsets enter these two checks.

The gap is also not an artifact of the onset detector. The four
model-agnostic baselines of Sec.~II-D bracket the identified moment
within a median $108$~mN\,m, and all detectors place the per-run onset
moment within $3$--$45$~mN\,m of one another, both well below
$\Delta M$.

\paragraph{Effect on the identified offset}
The reviewer's concern is whether this bias corrupts the delivered
offset, and the inversion answers it directly. Restoring both channels
in the onset balance and pairing the directions,
\begin{equation}
S_{\mathrm{off}} = \tfrac{1+\eta_M^{GE}}{2}(M_+ + M_-)
+ \tfrac{\eta_f^{GE}}{2}(f_+ l_+ - f_- l_-)
- \tfrac12\big[(W - f_+)l_+ - (W - f_-)l_-\big],
\label{eq:offge}
\end{equation}
the last bracket being dropped in the pivot-free variant of
\eqref{eq:pivotfree}, which assumes $l_+ = l_-$. Table~\ref{tab:offge}
evaluates \eqref{eq:offge} against the load-cell truth at both levels
of physics and both treatments of the arms.

\begin{table}[t]
\centering
\caption{Identified CoM offset against the load-cell truth over the
ten case/axis configurations, with and without the ground-effect
channels of \eqref{eq:offge}. ``$\Delta$'' is the mean paired change
in $|$error$|$ against the $\eta^{GE} = 0$ row of the same block, with
the exact sign test beside it.}
\label{tab:offge}
\small
\setlength{\tabcolsep}{4pt}
\begin{tabular}{@{}ll rrr rr@{}}
\hline
arms & ground effect & RMS & mean & max & $\Delta$ & $p$ \\
 & & [mm] & [mm] & [mm] & [mm] & \\
\hline
pivot-free & none (reported) & $1.64$ & $-0.91$ & $3.34$ & --- & --- \\
pivot-free & single-rotor & $1.67$ & $-0.98$ & $3.44$ & $-0.00$ & $1.00$ \\
pivot-free & interference & $1.82$ & $-1.20$ & $3.77$ & $+0.19$ & $0.11$ \\
\hline
pivot-based & none & $1.77$ & $+0.42$ & $2.96$ & --- & --- \\
pivot-based & single-rotor & $1.73$ & $+0.35$ & $2.80$ & $-0.01$ & $1.00$ \\
pivot-based & interference & $1.65$ & $+0.13$ & $3.05$ & $-0.06$ & $1.00$ \\
\hline
\end{tabular}
\end{table}

The ground-effect term does not improve the offset, and no paired
comparison reaches significance: the correction is $0.41$~mm in the
median against a $0.45$~mm standard error of the comparison itself, so
this dataset bounds the ground-effect contribution to the offset at
about half a millimetre either way rather than resolving it. What does
move systematically is the mean bias, which the fullest inversion ---
arms and interference together --- removes, from $-0.91$ to
$+0.13$~mm at unchanged RMS.

\paragraph{What is applied, and what is reported}
The two uses of the same measurement need separating. The compensation
applied at take-off is $M_{\mathrm{ff}}$ of \eqref{eq:pivotfree} with
no ground-effect model anywhere: it is the in-situ balanced moment and
already carries whatever the ground contributes, which is precisely
what has to be cancelled, so subtracting a modelled term there would
remove part of the disturbance the compensation exists to reject.
Reading the same quantity as a mass property is the separate question
of Table~\ref{tab:offge}, and the offset is reported at
$\eta^{GE} = 0$ to match the quantity flown, with the table standing
as the bound on what the omission costs.

That the deliverable is insensitive is structural rather than
fortunate. Writing the unmodelled part of the balance as a lumped
$\Delta_{\pm}$ and splitting it by direction,
$\Delta_a = (\Delta_+ - \Delta_-)/2$ and
$\Delta_s = (\Delta_+ + \Delta_-)/2$, everything acting as
force~$\times$~pivot-distance changes sign with the tip direction and
lands in $\Delta_a$, which cancels identically in
\eqref{eq:pivotfree}. Measured here, $|\Delta_a| = 54$~mN\,m in the
median against $|\Delta_s| = 37$, the latter $3.7\%$ of the half-width
of the interval $[M_-, M_+]$ the excitation measures directly
($1000$~mN\,m in the median). Since any moment strictly inside that
interval keeps both contacts loaded and the midpoint maximises the
smaller of the two margins, $M_{\mathrm{ff}}$ is the max--min-margin
choice whatever $\Delta_{\pm}$ turns out to be, and both endpoints
come from the excitation rather than from a model.

\paragraph{Limits of the attribution}
The comparison cannot assign $\Delta M$ uniquely. A contact-geometry
explanation fits equally well: if the normal-force resultant acts a
distance $\delta l$ inboard of the kinematic pivot circle, as a
finite compliant contact patch implies, a least-squares fit gives
$\delta l = 19.2 \pm 1.8$~mm with a residual indistinguishable from
that of the ground-effect hypothesis, and the same shift closes the
two checks above. The sensitivity that sets the floor is
$\mathrm{d}M/\mathrm{d}l = W - f = 10.6$~N, so the $61$~mN\,m the
interference model leaves is $5.8$~mm of contact arm --- the same
order as the run-to-run spread of the circle fit itself. Separating
the two requires a clearance sweep, since the aerodynamic term scales
with $h/R$ and the contact term does not: repeating one axis at
$h = 0.60$~m would put the aerodynamic hypothesis at $0.069$~N\,m
against the contact hypothesis's unchanged $0.20$, a three-fold split
well above the $0.06$~N\,m floor.
